% ==========================================================
% titles/title-06.tex -- Title VI: Economic Structure
% ==========================================================

\ConstTitle{Economic Structure}{title:economic}

\Article{The Three-Sector Economy}{art:three-sector-economy}

\ArtSection{The Three-Sector Structure}{art:three-sector-economy:s-three-sector-structure}
The Commonwealth economy is organized into three sectors, each governed according to the nature of the ownership relationship it involves:
\begin{itemize}
  \item The Worker Cooperative Sector, comprising most enterprises producing goods and services
  \item The Commons Sector, comprising natural resources, essential infrastructure, and shared services
  \item The Personal Sector, comprising individual consumption goods and small personal production
\end{itemize}

No sector may be abolished or restructured in ways that create private ownership of the means of production, by any legislative act, emergency declaration, or amendment of any kind, including under Article \ref{art:ordinary-amend}. Such restructuring constitutes arbitrary power over the survival and self-determination of others within the meaning of Article \ref{art:antidom}, \ref{art:antidom:s-ownership-domination}, and is therefore foreclosed by the Anti-Domination Principle's status as a hard limit under Article \ref{art:const-assembly}, \ref{art:hardlimits}.

\ArtSection{The Classification Test}{art:three-sector-economy:s-classification-test}
An economic activity is classified as a Commons asset under Article \ref{art:commons-sector}, \ref{art:commons-sector:s-commons-assets}, rather than organized within the Worker Cooperative Sector under Article \ref{art:worker-coop-sector}, where it satisfies either of the following:
\begin{itemize}
  \item \textbf{Natural monopoly or finite shared resource} -- the activity involves a physically finite resource, or infrastructure whose competitive duplication would be wasteful or physically incoherent, such that unified stewardship serves residents better than competitive provision. Land, water, energy infrastructure, telecommunications infrastructure, transport infrastructure, the financial system, and natural resources satisfy this test.
  \item \textbf{Care labor under conditions of impaired consumer choice} -- the activity constitutes care labor of the kind identified in Article \ref{art:equality}, \ref{art:equality:s-gender-equality}, delivered under conditions where the recipient cannot meaningfully evaluate quality or exercise informed choice at the point of need, such that market discipline structurally cannot function as a quality safeguard. Healthcare, education, childcare, and elder care satisfy this test.
\end{itemize}

An activity satisfying neither test -- including food production, construction, and general goods and services -- is organized within the Worker Cooperative Sector under Article \ref{art:worker-coop-sector}, subject to the essential-sector provisions of Article \ref{art:essential-labor} where applicable.

Where a new or emerging sector's classification is contested, the \gls{gls:sortition-legislature} applies this test by simple majority, subject to \gls{gls:constitutional-court} review confirming the classification is consistent with this Section. Reclassification of an activity from Commons to cooperative organization, or from cooperative to Commons, follows the same process and is additionally subject to \ref{art:three-sector-economy:s-three-sector-structure}'s prohibition on any revision that creates private ownership of the means of production.

\ArtSection{The Market Coordination Acknowledgment}{art:three-sector-economy:s-market-coordination-acknowledgment}
The Commonwealth retains market coordination among worker cooperatives as the primary mechanism for allocating most goods and services. This choice is made not because market coordination is consistent with the full elimination of domination, but because the material conditions of a complex modern economy have not produced a demonstrated alternative capable of coordinating production and distribution at scale without generating worse forms of concentrated power. Central planning at national scale has historically concentrated decision-making authority in ways that reproduce domination in new institutional forms. The cooperative market system is therefore a pragmatic response to current material conditions, not a final settlement.

The Commonwealth acknowledges that market competition among cooperatives generates its own structural pressures -- toward cost reduction, productivity intensification, and the gradual accumulation of competitive advantage by larger or better-positioned federations -- that the wage ratio cap, the antitrust provisions, and the Commons sector are designed to constrain but cannot fully eliminate. These contradictions are the primary object of ongoing constitutional self-analysis.

Future Constitutional Assemblies are explicitly invited to revisit this choice as material conditions change, as planning technology develops, and as the cooperative sector generates empirical evidence about whether market coordination can be transcended without reproducing the concentration of power it replaces. The three-sector structure may be revised under Article \ref{art:ordinary-amend}, subject to the constraint that no revision may create private ownership of the means of production or diminish the Material Sufficiency Floor.

\Article{The Worker Cooperative Sector}{art:worker-coop-sector}

\ArtSection{The Cooperative Mandate}{art:worker-coop-sector:s-cooperative-mandate}
Enterprises above a minimum size threshold -- set by the \gls{gls:national-cooperative-standards-body}, constituted under Article \ref{art:regulatory-bodies}, \ref{art:regulatory-bodies:s-major-regulatory-domains} -- must be organized as worker cooperatives in which all workers are member-owners with equal governance rights. For purposes of determining whether the minimum size threshold is met, the \gls{gls:national-cooperative-standards-body} looks through arrangements where multiple nominally separate enterprises are commonly controlled, contractually bound to operate as a single economic unit, or structured with the primary purpose of remaining under the threshold, and aggregates them accordingly.

\ArtSection{Ownership Structure}{art:worker-coop-sector:s-ownership-structure}
Cooperative ownership stakes may not be sold to non-workers, transferred to investors, or inherited. Workers receive stakes automatically after a qualifying period determined by the cooperative's democratic governance. Upon departure, workers receive fair compensation for accumulated contribution; the stake itself reverts to the cooperative.

Civic integration status does not affect a worker's ownership rights, governance rights, or surplus distribution rights within their cooperative. Worker-membership in a cooperative, including full ownership rights and democratic governance participation within that cooperative, is available to all residents regardless of civic integration status. Cooperative ownership rights are governed by this Article and are distinct from political membership rights conferred by full civic integration under Article \ref{art:civic-integration}.

\ArtSection{Democratic Governance}{art:worker-coop-sector:s-democratic-governance}
Cooperative governance operates on the principle of one worker, one vote for all major decisions. Day-to-day management may be delegated to elected managers who are accountable to and recallable by worker-owners.

\ArtSection{The Wage Ratio Cap}{art:worker-coop-sector:s-wage-ratio-cap}
The total compensation -- including wages, surplus distributions, and benefits -- of the highest-paid worker in any cooperative may not exceed ten times the total compensation of the lowest-paid worker. This ratio may be adjusted downward but never upward by the \gls{gls:sortition-legislature}.

\ArtSection{Cooperative Markets}{art:worker-coop-sector:s-cooperative-markets}
Cooperatives compete with each other in markets for most goods and services. No cooperative may acquire a dominant market position that enables it to exercise power over other cooperatives or over consumers' access to essential goods. The \gls{gls:financial-regulatory-body}, constituted under Article \ref{art:regulatory-bodies}, \ref{art:regulatory-bodies:s-major-regulatory-domains}, is the enforcement authority for this prohibition. Affected cooperatives and consumers have standing to bring complaints directly to the Body.

\ArtSection{The Cooperative Development Bank}{art:worker-coop-sector:s-cooperative-development-bank}
A publicly chartered \gls{gls:cooperative-development-bank} provides startup financing and working capital to cooperatives at cost rather than for profit. The Bank is governed by a fifteen-member board, selected by sortition from the cooperative sector on staggered five-year non-renewable terms, and is accountable to the \gls{gls:sortition-legislature}.

The \gls{gls:cooperative-development-bank} may not develop, maintain, or use any credit scoring system that assesses the creditworthiness of cooperatives or their members against historical financial records for the purpose of determining lending eligibility. This prohibition applies to internally developed scoring systems and to any third-party scoring data. Lending eligibility is determined exclusively through the Development Partnership process under \ref{art:worker-coop-sector:s-development-partnership-process}.

The Bank maintains portfolio health monitoring of all financed cooperatives for the sole purpose of early identification of financial distress and activation of support relationships. Portfolio health data may not be used to restrict, reduce, or recall financing, and may not be shared with any external entity without the consent of the cooperative. The \gls{gls:citizen-oversight-panels} reviews portfolio monitoring practices annually and publishes findings through the \gls{gls:transparency-engine}.

\ArtSection{The Development Partnership Process}{art:worker-coop-sector:s-development-partnership-process}
The Bank assesses cooperative viability through a Development Partnership process that builds capacity alongside assessment. A Bank development team works with prospective cooperatives to refine the business model, identify capital needs, connect the cooperative to training and Commons infrastructure, and assess labor commitment. Lending decisions are made by rotating Peer Panels drawn from worker-members of established cooperatives in adjacent sectors, selected by sortition from a pool of volunteers, operating against published criteria. No application may be declined solely on the basis of absence of collateral, prior credit history, or members' personal financial history.

Essential sector cooperatives in underserved regions are eligible for Commons-backstop grants in addition to or in lieu of lending. The Bank maintains ongoing support relationships with all financed cooperatives, with distress signals triggering support rather than enforcement responses. At founding, Peer Panels are drawn from the \gls{gls:constitutional-assembly}'s designees, with explicit sunset once the cooperative sector reaches one hundred established enterprises. See Article \ref{art:founding-mechanisms}, \ref{art:founding-mechanisms:s-cooperative-development-bank}.

\ArtSection{The Cooperative Federation}{art:worker-coop-sector:s-cooperative-federation}
Voluntary federations of cooperatives in related sectors may share research, training, supply chains, and mutual insurance. Federation membership is voluntary. Federations may not develop governance structures that concentrate power over member cooperatives. Federations conduct sector-knowledge pre-assessments of prospective cooperatives before applications reach the \gls{gls:cooperative-development-bank}; this role is advisory and the Bank is not bound by federation assessments. Cooperatives in sectors without an established federation proceed directly to the Bank's Development Partnership process.

\Article{The Commons Sector}{art:commons-sector}

\ArtSection{Commons Assets}{art:commons-sector:s-commons-assets}
The following assets are held in common by all residents of the Commonwealth:
\begin{itemize}
  \item All land
  \item Water systems and water resources
  \item Energy generation infrastructure and distribution grids
  \item Telecommunications infrastructure
  \item Transportation infrastructure including roads, rail, ports, and airports
  \item The core financial system including payment infrastructure and the monetary authority
  \item Natural resources including minerals, forests, and fisheries
  \item The healthcare delivery system
  \item The education system at all levels
  \item The childcare and elder care systems
\end{itemize}

\ArtSection{General Principles}{art:commons-sector:s-general-principles}
All assets listed in \ref{art:commons-sector:s-commons-assets} are held in common permanently and may not be privatized in whole or in part by any legislative act, emergency declaration, or amendment of any kind, including under Article \ref{art:ordinary-amend}. Privatization of Commons assets constitutes arbitrary power over the survival of others within the meaning of Article \ref{art:antidom}, \ref{art:antidom:s-ownership-domination}, and is therefore foreclosed by the Anti-Domination Principle's status as a hard limit under Article \ref{art:const-assembly}, \ref{art:hardlimits}.

Commons assets are managed exclusively to maximize use value for all residents; surplus generated flows to the Material Sufficiency Floor and Commons reinvestment under Article \ref{art:land-tenures}, \ref{art:land-tenures:s-operating-principle}. Access to Commons assets may not be rationed on the basis of ability to pay, civic integration status, or any characteristic protected under Article \ref{art:equality}. All Commons asset governance is fully transparent through the \gls{gls:transparency-engine} under Article \ref{art:transparency-engine}.

\ArtSection{Water}{art:commons-sector:s-water}
Water is a biological necessity and its governance is a foundational obligation of the Commonwealth. Water systems and water resources are managed for universal access, ecological health, and intergenerational sustainability. Extraction, diversion, or degradation of water resources is subject to mandatory ecological impact assessment under Article \ref{art:ecological-limits}. Pricing or rationing of water on the basis of ability to pay is prohibited. The rights of indigenous nations over water resources on or adjacent to returned land are determined through the Land Back negotiation process under Article \ref{art:indigenous-sov}, and Commonwealth water governance does not supersede those rights.

\ArtSection{Energy}{art:commons-sector:s-energy}
Energy infrastructure is managed to provide universal access and to fulfill the Commonwealth's obligations under the climate mandate of Article \ref{art:ecological-limits}. Energy pricing may not create energy poverty -- defined as a household spending more than a defined proportion of its income on energy, set by the relevant \gls{gls:implementation-councils} and subject to \gls{gls:sortition-legislature} review. The transition of energy generation from extraction-based to renewable systems is a constitutional obligation whose timeline and escalation triggers are governed by Article \ref{art:ecological-limits}. No new extraction-based energy infrastructure may be approved that is inconsistent with the current climate mandate targets.

\ArtSection{Telecommunications}{art:commons-sector:s-telecommunications}
Telecommunications infrastructure is managed to provide universal access as a precondition of civic participation, education, and economic activity. Access may not be rationed on the basis of ability to pay or any characteristic protected under Article \ref{art:equality}. The telecommunications infrastructure may not be used as an instrument of content surveillance, behavioral monitoring, or population tracking. Mass collection of communications or behavioral data through telecommunications infrastructure is prohibited and governed by Article \ref{art:internal-security}, \ref{art:internal-security:s-no-mass-surveillance}. Network access is provided on equal terms to all users; no prioritization of content, source, or user class is permitted.

\ArtSection{Transport}{art:commons-sector:s-transport}
Transport infrastructure is managed for universal access with particular attention to regional equity, ensuring that rural and underserved communities are not systematically disadvantaged in access to Commonwealth services, economic activity, and civic participation. All significant transport infrastructure decisions are subject to ecological impact assessment under Article \ref{art:ecological-limits}.

\Article{Land and Tenures}{art:land-tenures}

\ArtSection{Residential Land}{art:land-tenures:s-residential-land}
Persons occupy residential land through one of two mechanisms: Commons-Allocated Housing under Article \ref{art:open-presence}, \ref{art:open-presence:s-commons-allocated-housing}, which is the primary mechanism through which all persons access housing; or a Residential Tenure, recognized where a person has constructed a dwelling or made improvements whose assessed labor value exceeds twenty-five percent of the assessed replacement cost of the dwelling at the time of the claim. Residential Tenures are assessed by the \gls{gls:housing-implementation-council} against published criteria.

No person may hold more than one Residential Tenure acquired through their own improvements simultaneously. Where a person who holds a Residential Tenure makes improvements to a second dwelling that would otherwise qualify for a Residential Tenure, they must elect which dwelling retains the Tenure. Making this election does not affect the person's right to occupy either dwelling; it determines which dwelling carries heritable tenure status. The non-elected dwelling is held as Commons-Allocated Housing -- the person may continue to occupy it, but it returns to Commons housing stock rather than passing to their heirs upon their death or departure. The second dwelling is treated as Commons-Allocated Housing from the moment the improvements are made until the election is made, and as Commons-Allocated Housing thereafter if not elected.

Residential Tenures are secure, heritable within immediate family, and non-transferable on any market. They may be used as collateral for lending exclusively through the \gls{gls:public-banking-system} under Article \ref{art:finance-money}, \ref{art:finance-money:s-public-banking}. Default on a Residential Tenure used as collateral may not result in loss of that Tenure; the \gls{gls:public-banking-system}'s remedy in default is limited to a claim on improvement compensation proceeds at the point of return to Commons stock. Upon the holder's death, the Tenure passes to immediate family members who occupy the dwelling as their primary residence. Where no inheriting party occupies the dwelling as their primary residence for a continuous period exceeding three years, the Tenure expires and the dwelling returns to Commons housing stock. Fair compensation for improvements is paid to the estate at assessed replacement cost. Temporary non-occupancy due to medical incapacity, essential service obligations, or other circumstances recognized by the \gls{gls:housing-implementation-council} against published criteria does not count toward the three-year threshold.

\ArtSection{Inherited Tenures}{art:land-tenures:s-inherited-tenures}
A Residential Tenure received through inheritance is held as an Inherited Tenure, distinct from any Residential Tenure acquired through the holder's own improvements. Where inheritance would result in a person holding more than one Inherited Tenure, the heirs determine which Tenure to retain and which to return to Commons stock with fair compensation paid to the estate. No person may hold more than two Tenures in total at any time, regardless of how they were acquired. The three-year non-occupancy clock is suspended on all affected properties until the heirs have made their determination. Upon determination, the retained Tenure is subject to normal non-occupancy provisions; properties returned to Commons stock return immediately with improvement compensation paid to the estate.

\ArtSection{Productive Tenures}{art:land-tenures:s-productive-tenures}
Cooperatives may hold Productive Tenures over land used for cooperative operations. Productive Tenures are distinct from Residential Tenures in all respects and are governed by this Section and Article \ref{art:worker-coop-sector}.

Cooperatives access Productive Tenures through one of two mechanisms: Commons allocation administered by the relevant \gls{gls:implementation-councils} for the land's domain, which is the primary mechanism; or a Builder's Productive Tenure, recognized where a cooperative has constructed its premises or made improvements whose assessed labor value exceeds fifty percent of the assessed replacement cost at the time of the claim, assessed by the relevant \gls{gls:implementation-councils} against published criteria.

Productive Tenures are secure and non-transferable on any market. When a cooperative dissolves, relocates, or permanently vacates its premises, the tenure returns to Commons stock. The cooperative receives fair compensation for improvements at assessed replacement cost. Where a dissolving cooperative wishes its premises to pass to a successor cooperative in the same or adjacent sector, a managed non-market transfer may be approved by the relevant \gls{gls:implementation-councils}. No price is set by market negotiation in any transfer; compensation is assessed at replacement cost of improvements only.

Productive Tenures may be used as collateral for lending exclusively through the \gls{gls:cooperative-development-bank} under Article \ref{art:worker-coop-sector}, \ref{art:worker-coop-sector:s-cooperative-development-bank}, or the \gls{gls:public-banking-system} under Article \ref{art:finance-money}, \ref{art:finance-money:s-public-banking}. Default on a Productive Tenure used as collateral may not result in loss of that tenure; the lending institution's remedy in default is limited to a claim on improvement compensation proceeds at the point of return to Commons stock.

\ArtSection{Acknowledgment}{art:land-tenures:s-acknowledgment}
The Commonwealth acknowledges that the Commons allocation system, as the primary mechanism for housing access, cannot perfectly replicate the granular individual choice a housing market provides. Persons seeking a specific dwelling they did not build receive Commons-Allocated Housing governed by objective criteria rather than market selection. This constraint is accepted because the alternative -- a market in housing -- reproduces domination by conditioning access to a fundamental right on financial resources. Future Constitutional Assemblies are explicitly invited to revisit the allocation mechanism as Commons housing stock matures, allocation systems develop, and empirical evidence accumulates about whether more responsive allocation is achievable without reproducing market domination.

\ArtSection{Land Value and Displacement Protection}{art:land-tenures:s-land-value-displacement}
Land value generated by surrounding community development, infrastructure investment, and collective social activity may not be privately captured. Since all tenure forms recognized under this Article are non-transferable on any market, appreciation in land value remains with the Commons and is realized upon reversion of any tenure to Commons stock. The Commons Land Fund, which supports housing infrastructure investment and the Material Sufficiency Floor, is funded through Commons sector operating surplus and \gls{gls:monetary-authority} allocation under Article \ref{art:finance-money}, \ref{art:finance-money:s-commons-investment-fund}.

\ArtSection{Displacement Protection}{art:land-tenures:s-displacement-protection}
No person occupying a dwelling through Commons-Allocated Housing or a Residential Tenure may be displaced without a public interest determination by a regional Commons body, provision of equivalent replacement housing prior to displacement, full compensation for improvements at assessed replacement cost for Residential Tenure holders, and access to \gls{gls:constitutional-court} review. This protection applies without exception to every person on Commonwealth territory regardless of occupancy status.

\ArtSection{Commons Governance}{art:land-tenures:s-commons-governance}
Commons assets are governed at the level appropriate to their scope. Local Commons are governed by local community assemblies. Regional Commons are governed by regional assemblies. National Commons are governed by the \gls{gls:sortition-legislature} with implementation by specialized Commons management bodies. All Commons governance is fully transparent and subject to public participation.

\ArtSection{Operating Principle}{art:land-tenures:s-operating-principle}
Commons entities are managed to maximize use value for their communities, not to generate surplus for extraction. Where surplus is generated, it flows directly to the Material Sufficiency Floor and Commons reinvestment.

\Article{Finance and Money}{art:finance-money}

\ArtSection{Prohibited Financial Activities}{art:finance-money:s-prohibited-financial-activities}
The following financial activities are prohibited as mechanisms of rent extraction inconsistent with the anti-domination principle:
\begin{itemize}
  \item Speculative derivative instruments whose primary function is extraction of value without production
  \item Private equity acquisition of cooperatives for asset stripping or resale
  \item Any financial instrument that creates shareholder extraction from worker cooperative surplus
  \item High-frequency trading designed to exploit information asymmetries
  \item Private lending -- the provision of credit by any non-Commons entity in exchange for interest, fees, or any other form of return.
\end{itemize}

All lending functions are performed exclusively by Commons institutions: the \gls{gls:public-banking-system} for personal lending, the \gls{gls:cooperative-development-bank} for cooperative lending, and the \gls{gls:commons-investment-fund} for infrastructure investment. No person, cooperative, or other entity may extend credit for return outside these institutions. Informal lending between individuals without interest or fees is not governed by this prohibition. Residential Tenure lending is available exclusively through the \gls{gls:public-banking-system} under \ref{art:finance-money:s-public-banking}. Productive Tenure lending is available through the \gls{gls:public-banking-system} or the \gls{gls:cooperative-development-bank} under Article \ref{art:worker-coop-sector}, \ref{art:worker-coop-sector:s-cooperative-development-bank}. No tenure of either type may be used as collateral with any entity other than these two Commons lending institutions.

\ArtSection{Public Banking}{art:finance-money:s-public-banking}
A \gls{gls:public-banking-system} provides savings, checking, and personal lending services to all residents at cost. The system is a Commons institution governed by the \gls{gls:sortition-legislature}.

\ArtSection{The Monetary Authority}{art:finance-money:s-monetary-authority}
Money creation is a public function of the democratically governed \gls{gls:monetary-authority}, accountable to the \gls{gls:sortition-legislature}. Private money creation through fractional reserve lending is prohibited.

The \gls{gls:monetary-authority}'s allocation to Commons operations and the Material Sufficiency Floor is calibrated against real productive capacity in the Commonwealth economy. Monetary allocation is a stabilizing mechanism for timing mismatches between Commons revenue and Floor obligations; it is not a substitute for adequate cooperative surplus contributions. The \gls{gls:monetary-authority} publishes quarterly assessments of the relationship between money creation, productive capacity, and price stability through the \gls{gls:transparency-engine}. Where the \gls{gls:monetary-authority} determines that allocation is substituting for structurally inadequate contribution rates rather than smoothing timing variation, it is constitutionally obligated to report this finding to the \gls{gls:sortition-legislature} within thirty days. The Legislature must convene a mandatory review session within sixty days of receiving such a report.

A structural-dependence finding triggers an automatic minimum increase in cooperative surplus contribution rates, set jointly by the \gls{gls:monetary-authority} and the \gls{gls:financial-regulatory-body} at the level necessary to close the identified structural gap. This automatic increase takes effect at the mandatory review session unless the \gls{gls:sortition-legislature} affirmatively rejects it by the same supermajority required to lower contribution rates under \ref{art:finance-money:s-cooperative-surplus-contributions}, and publishes its reasoning through the \gls{gls:transparency-engine}. Rejection does not suspend future structural-dependence findings or future automatic increases arising from them.

\ArtSection{The Commons Investment Fund}{art:finance-money:s-commons-investment-fund}
Long-term infrastructure investment is financed through the \gls{gls:commons-investment-fund}, funded by cooperative surplus contributions, Commons sector operating surplus, and \gls{gls:monetary-authority} allocation. These fund the Commons Land Fund under Article \ref{art:commons-sector}, \ref{art:commons-sector:s-general-principles}. All allocations are fully transparent.

\ArtSection{Enforcement}{art:finance-money:s-enforcement}
The \gls{gls:financial-regulatory-body}, constituted under Article \ref{art:regulatory-bodies}, is the enforcement authority for the prohibitions in this Article. Investigation findings are submitted to the \gls{gls:security-enforcement-review-panel} for authorization of specific enforcement actions under Article \ref{art:internal-security}, \ref{art:internal-security:s-security-enforcement-review}. Financial violations with criminal dimensions are referred to the relevant justice \gls{gls:implementation-councils}. The \gls{gls:monetary-authority} is required to report to the \gls{gls:financial-regulatory-body} any financial flows it identifies as potentially inconsistent with this Article; this reporting obligation does not give the \gls{gls:monetary-authority} independent enforcement authority.

\ArtSection{Cooperative Surplus Contributions}{art:finance-money:s-cooperative-surplus-contributions}
The \gls{gls:sortition-legislature} sets cooperative surplus contribution rates at levels demonstrably sufficient to fund the Material Sufficiency Floor obligations established in Article \ref{art:matsuff}, together with Commons sector operating surplus and other Commons revenue streams. Contribution rates may not be set below the level the \gls{gls:monetary-authority} and the \gls{gls:financial-regulatory-body} jointly assess as necessary to fund Floor obligations without structural dependence on monetary allocation.

The \gls{gls:transparency-engine} monitors in real time the gap between Material Sufficiency Floor obligations and Commons revenue from all sources, including cooperative surplus contributions, Commons sector operating surplus, ecological resource revenues, and \gls{gls:monetary-authority} allocation. Where the monitored gap indicates that Floor obligations are not fully funded on a structural basis, a public alert is issued automatically and the Legislature must convene a mandatory review session within sixty days.

Contribution rates are reviewed every five years alongside the climate mandate under Article \ref{art:ecological-limits}. They may be raised by simple legislative majority and may be lowered only by supermajority vote with a published funding impact assessment confirmed by the \gls{gls:financial-regulatory-body}. No reduction in contribution rates may take effect if the \gls{gls:monetary-authority} has issued a structural dependence finding under \ref{art:finance-money:s-monetary-authority} within the preceding two years.

\ArtSection{Floor Growth Allocation}{art:finance-money:s-floor-growth-allocation}
Meeting Material Sufficiency Floor obligations satisfies the minimum requirement of this Article but does not exhaust it. No less than fifteen percent of the annual real growth in cooperative surplus contributions and Commons sector operating surplus, measured year-over-year and confirmed by the \gls{gls:financial-regulatory-body}, must be directed to expanding the scope, quality, or standard of Material Sufficiency Floor provision, beyond whatever level current obligations already require. This requirement may be satisfied only from real growth in cooperative surplus contributions and Commons sector operating surplus as so measured; \gls{gls:monetary-authority} allocation, and any revenue collected under \ref{art:finance-money:s-public-banking}, may never be counted toward it. In a year without qualifying real growth, no allocation is required. The Legislature publishes, through the \gls{gls:transparency-engine}, an annual accounting of qualifying growth and the specific Floor expansions it funded. Persistent qualifying growth with no corresponding Floor expansion over three consecutive years makes that year's accounting subject to Legislative Decision Review under Article \ref{art:civic-initiative-infrastructure}, \ref{art:civic-initiative-infrastructure:s-constitutional-intervention-types}, triggering the mandatory reconsideration session that intervention type provides.

\ArtSection{Personal Sector Wealth Transfer Contribution}{art:finance-money:s-wealth-transfer-contribution}
Productive and Residential Tenures are governed exclusively by the inheritance provisions of Article \ref{art:land-tenures:s-inherited-tenures} and are not affected by this Section. A wealth transfer contribution applies only to Personal Sector savings, personal property, and other non-tenure assets passing by inheritance or bequest. The contribution phases in linearly on the portion of a bequest to any single heir that exceeds five times the Commonwealth median annual wage, reaching a maximum rate on the portion exceeding twenty times the Commonwealth median annual wage, at rates set by the \gls{gls:sortition-legislature} and reviewed on the same five-year cycle as cooperative surplus contribution rates under \ref{art:finance-money:s-cooperative-surplus-contributions}. Bequests below five times the Commonwealth median annual wage are exempt entirely, and ordinary personal property of primarily sentimental or functional rather than financial value is exempt regardless of value. Revenue is directed to the Material Sufficiency Floor and logged through the \gls{gls:transparency-engine}.

\Article{Intellectual Property and Innovation}{art:ip-innovation}

\ArtSection{Abolition of Perpetual Rent}{art:ip-innovation:s-abolition-perpetual-rent}
Patents, perpetual copyright, and other instruments of permanent rent extraction from intellectual works are abolished. Inventors, creators, and researchers are entitled to fair compensation for their labor of creation, not permanent monopoly rents from others' use of their creations. The right of attribution -- to all persons involved in the invention, creation, research, or development of an intellectual work -- is permanent and may not be waived or transferred.

\ArtSection{The Knowledge Commons}{art:ip-innovation:s-knowledge-commons}
All research outputs from publicly funded institutions immediately become public domain. Knowledge produced within the Commonwealth is treated as a Commons asset.

\ArtSection{Innovation Incentives}{art:ip-innovation:s-innovation-incentives}
Innovation is incentivized through public research funding, innovation prizes for solutions to defined public challenges, and the intrinsic cooperative self-interest of worker-owners. Prize-winning innovations become public domain upon award.

\ArtSection{Creator Protections and Principled Duration}{art:ip-innovation:s-creator-protections-principled}
Creators of artistic and literary works receive protection against unauthorized reproduction. The duration of this protection is tied to the compensation structure:
\begin{itemize}
  \item Works for which the creator has accepted a public innovation prize or Commons commission enter public domain immediately upon award.
  \item Works produced independently receive protection for a period sufficient to allow genuine market compensation recovery. The outer protection limit is set at the point at which empirical revenue data for the relevant creative domain shows median annual revenue has declined below five percent of peak -- currently assessed at fifteen years from first publication, reflecting empirical revenue curves rather than an arbitrary number. This assessment is subject to revision by the \gls{gls:sortition-legislature} based on updated empirical data.
  \item Moral rights -- the right of attribution and the right of integrity -- are permanent and may not be waived or transferred.
  \item At the outer limit, all works enter the Knowledge Commons unconditionally.
\end{itemize}

\Article{Ecological Limits}{art:ecological-limits}

\ArtSection{Economic Purpose}{art:ecological-limits:s-economic-purpose}
Economic growth is not a goal of the Commonwealth's economic system. The goal is material sufficiency for all residents within ecological limits. Steady-state sufficiency takes precedence over expansion.

\ArtSection{Rights of Natural Systems}{art:ecological-limits:s-rights-natural-systems}
Natural ecosystems possess rights to exist, regenerate, and flourish. These rights are justiciable before the \gls{gls:constitutional-court}. Commons governance of land, water, and natural resources is conducted as stewardship, not extraction.

\ArtSection{Ecological Impact Assessment}{art:ecological-limits:s-ecological-impact-assessment}
All economic planning, Commons investment, and significant cooperative activity is subject to mandatory ecological impact assessment. Hard limits on ecosystem disruption may not be overridden by democratic majorities.

\ArtSection{Climate Mandate}{art:ecological-limits:s-climate-mandate}
The Commonwealth operates under a rolling climate mandate reviewed by the \gls{gls:sortition-legislature} every five years, with automatic escalation triggers tied to objective scientific metrics.

\ArtSection{Emissions Fee}{art:ecological-limits:s-emissions-fee}
In addition to the emissions standards enforced by the \gls{gls:environmental-regulatory-body} under \ref{art:regulatory-bodies:s-major-regulatory-domains}, an emissions fee is levied on Worker Cooperative Sector and Commons Sector production, priced by the \gls{gls:environmental-regulatory-body} against the same objective scientific metrics that govern the climate mandate, and reviewed on the same five-year cycle with the same automatic escalation triggers. The fee applies only to activity within the hard limits set under \ref{art:ecological-limits:s-ecological-impact-assessment}; it has no authority to raise, waive, or substitute for those hard limits, which remain enforced directly regardless of any resident's or cooperative's ability or willingness to pay. The fee does not apply to Personal Sector consumption. Revenue is divided between the \gls{gls:commons-investment-fund}, prioritizing investment that reduces future emissions, and the Material Sufficiency Floor, in a published split counted as ecological resource revenue under \ref{art:finance-money:s-cooperative-surplus-contributions} and logged in full through the \gls{gls:transparency-engine}.

\Article{Essential Labor and Material Production}{art:essential-labor}

\ArtSection{The Foundational Obligation}{art:essential-labor:s-foundational-obligation}
The Commonwealth acknowledges that the material rights guaranteed in Article \ref{art:matsuff} depend on sufficient human labor to produce food, deliver healthcare, teach students, build housing, and maintain infrastructure. The system is designed to make this labor attractive rather than coerced.

\ArtSection{Essential Labor}{art:essential-labor:s-essential-labor}
Essential labor is organized either through cooperatives under Article \ref{art:worker-coop-sector} or directly through Commons institutions under Article \ref{art:commons-sector}, depending on the sector's classification under Article \ref{art:three-sector-economy}, \ref{art:three-sector-economy:s-classification-test}.

\textbf{Cooperative-organized essential sectors.} Cooperatives in essential sectors -- food production, construction, and other essential goods and services organized under Article \ref{art:worker-coop-sector} -- are subject to constitutional minimum compensation floors. The lowest-paid worker in any essential sector cooperative may not earn below a defined ratio of the overall cooperative sector wage median, set by the \gls{gls:sortition-legislature} and may be raised but not lowered.

\textbf{Commons institution essential labor.} Workers in healthcare, education, childcare, and elder care are directly employed by Commons institutions under Article \ref{art:commons-sector} rather than organized as worker-owners, and receive no surplus distribution or ownership stake. The lowest-paid worker in any such Commons institution may not earn below a compensation floor set by the \gls{gls:sortition-legislature}, calibrated to remain financially equivalent to that of an essential-sector cooperative worker-owner unless the Legislature determines a different floor is warranted for a specific role or institution. This floor may be raised but not lowered. The relevant \gls{gls:implementation-councils} for each Commons institution sets an internal compensation structure consistent with this floor; compensation structures are published and logged through the \gls{gls:transparency-engine}, and reviewed annually by the \gls{gls:citizen-oversight-panels} embedded within the \gls{gls:implementation-councils} under Article \ref{art:citizen-oversight-panels}.

No coercive assignment of persons to labor under either subsection is permissible under any circumstances.

\ArtSection{Training Pipelines}{art:essential-labor:s-training-pipelines}
Training and education pipelines for all essential vocations are Commons institutions, fully funded, regionally distributed, and free to all who enter them. Geographic distribution is designed to address the concentration problem that depletes essential services in rural and underserved areas.

Persons enrolled in essential sector training programs receive a training stipend for the duration of their enrolled program. The stipend is set at no less than the applicable essential sector compensation floor established in \ref{art:essential-labor:s-essential-labor}, ensuring that entering a training pipeline is financially equivalent to entry-level essential sector employment. The stipend may be raised above this floor by the \gls{gls:sortition-legislature}, including in response to pipeline monitoring triggers under \ref{art:essential-labor:s-essential-sector-pipeline}. The stipend may not be reduced below the \ref{art:essential-labor:s-essential-labor} floor by any legislative act or \gls{gls:implementation-councils} decision. Stipends are administered by the relevant Commons training institution and logged through the \gls{gls:transparency-engine}.

\ArtSection{Voluntary Service Incentives and Automatic Escalation}{art:essential-labor:s-voluntary-service-incentives}
Persons who complete defined service periods in underserved essential roles receive enhanced Commons benefits including priority weighting in Commons-Allocated Housing under Article \ref{art:open-presence}, \ref{art:open-presence:s-stable-housing-transition}, extended sabbatical rights, and enhanced retirement security. This service is voluntary, subject to the prohibition on coercive assignment established in \ref{art:essential-labor:s-essential-labor}.

The relevant \gls{gls:implementation-councils} maintains a published Incentive Escalation Schedule, tiered against the severity of a shortfall determined under \ref{art:essential-labor:s-essential-sector-pipeline} or \ref{art:essential-labor:s-forward-looking-capacity}, establishing in advance which specific enhanced benefits activate at each tier of severity. Later tiers may not offer fewer benefits than earlier tiers.

Upon a forward-looking shortfall identified under \ref{art:essential-labor:s-forward-looking-capacity}, or a reactive threshold crossing under \ref{art:essential-labor:s-essential-sector-pipeline}, escalation to the corresponding tier is automatic and immediate, activating at the same time as the mandatory legislative review session those Sections require, rather than awaiting that session's outcome. The \gls{gls:sortition-legislature} retains authority to expand incentives beyond the published schedule at its review session, and may revise the Escalation Schedule itself for future shortfalls, but may not suspend or reduce benefits already activated for persons currently serving.

Escalating benefits may include additional individual incentives as determined by the \gls{gls:sortition-legislature}, subject to the wage ratio cap of Article \ref{art:worker-coop-sector}, \ref{art:worker-coop-sector:s-wage-ratio-cap}, and may separately include enterprise- or federation-level capacity investment -- including but not limited to expanded Commons-backstop grants under Article \ref{art:worker-coop-sector}, \ref{art:worker-coop-sector:s-development-partnership-process}, preferential \gls{gls:cooperative-development-bank} financing terms, or priority Productive Tenure allocation -- which is not compensation to any individual worker and is not subject to that cap. Escalating benefits may never take the form of coercive assignment of any kind.

Escalation Schedules are published and logged through the \gls{gls:transparency-engine} and reviewed every five years alongside the climate mandate under Article \ref{art:ecological-limits}.

\ArtSection{Automation of Burdensome Labor}{art:essential-labor:s-automation-burdensome-labor}
The Commonwealth actively pursues automation of the most physically dangerous, repetitive, or degrading elements of essential work. Efficiency gains from automation in essential sectors flow to the Commons rather than to cooperative surplus, and should reduce the labor burden instead of displacing workers.

\ArtSection{Essential Sector Pipeline Monitoring}{art:essential-labor:s-essential-sector-pipeline}
Each \gls{gls:implementation-councils} with responsibility for an essential sector maintains a continuous public register of student enrollment in training pipelines for that sector, projected workforce needs over a ten-year horizon, and current workforce capacity. Methodology for projections and threshold determinations is published in full through the \gls{gls:transparency-engine} and reviewed every five years alongside the climate mandate under Article \ref{art:ecological-limits}.

When enrollment in any essential sector pipeline falls below a threshold determined by the relevant \gls{gls:implementation-councils} to indicate risk of future workforce insufficiency, the Council is constitutionally obligated to report immediately to the \gls{gls:coordination-assembly} and to the \gls{gls:sortition-legislature}. The Legislature must convene a mandatory review session within sixty days to assess and strengthen voluntary incentive mechanisms under \ref{art:essential-labor:s-voluntary-service-incentives}.

\gls{gls:implementation-councils}s report pipeline status annually to the \gls{gls:coordination-assembly}. The Assembly identifies cross-sector patterns and refers systemic findings to the Legislature with recommended responses. This remains subject to the prohibition on coercive assignment established in \ref{art:essential-labor:s-essential-labor}, even when pipeline thresholds are crossed.

\ArtSection{Forward-Looking Capacity Targets}{art:essential-labor:s-forward-looking-capacity}
In addition to the threshold monitoring established in \ref{art:essential-labor:s-essential-sector-pipeline}, each \gls{gls:implementation-councils} with responsibility for an essential sector sets forward-looking workforce capacity targets, projected over the same ten-year horizon, rather than relying solely on reactive enrollment monitoring. Targets are derived from projected demand for the sector's output -- including, where relevant, projected population growth, immigration trends under Article \ref{art:open-presence}, and any housing gap identified under Article \ref{art:open-presence}, \ref{art:open-presence:s-stable-housing-transition} -- rather than from current enrollment or current workforce levels alone.

For the construction sector specifically, capacity targets are calibrated against the \gls{gls:housing-implementation-council}'s housing gap projections under Article \ref{art:open-presence}, \ref{art:open-presence:s-stable-housing-transition}, and are set at a level intended to prevent a housing gap from arising, not merely to close one once identified.

Where projected capacity is forecast to fall short of a forward-looking target at any point within the ten-year horizon, the responsible \gls{gls:implementation-councils} is constitutionally obligated to report this finding to the \gls{gls:coordination-assembly} and the \gls{gls:sortition-legislature} immediately upon the shortfall's identification, without waiting for the reactive enrollment threshold of \ref{art:essential-labor:s-essential-sector-pipeline} to be crossed. The Legislature must convene a mandatory review session within sixty days to strengthen voluntary incentive mechanisms under \ref{art:essential-labor:s-voluntary-service-incentives}, expand training pipeline capacity under \ref{art:essential-labor:s-training-pipelines}, or both.

Forward-looking targets and their underlying methodology are published in full through the \gls{gls:transparency-engine} and reviewed every five years alongside the climate mandate under Article \ref{art:ecological-limits}, consistent with \ref{art:essential-labor:s-essential-sector-pipeline}. This remains subject to the prohibition on coercive assignment established in \ref{art:essential-labor:s-essential-labor}, even where forward-looking targets are unmet.

\ArtSection{Food Stability Implementation}{art:essential-labor:s-food-stability-implementation}
Food production remains within the Worker Cooperative Sector under Article \ref{art:worker-coop-sector} and is not converted to direct Commons provision. The food security right guaranteed under Article \ref{art:matsuff} is instead delivered as an unconditional entitlement to a defined basket of staple foods, redeemable by any resident at any food-producing or food-retailing cooperative. Redeeming cooperatives are reimbursed at cost by the Commons through the mechanism established in \ref{art:finance-money:s-cooperative-surplus-contributions}. The entitlement is not means-tested, is not reduced by a resident's other income or purchases, and is not affected by price movement in the non-entitlement portion of the food market. The composition of the staple basket is set by the relevant \gls{gls:implementation-councils}, published through the \gls{gls:transparency-engine}, and may be expanded but not narrowed. Hospitality, prepared food, and dining services are not part of this entitlement and are ordinary Worker Cooperative Sector economic activity governed by \ref{art:worker-coop-sector} in full.

\ArtSection{Essential Goods Market Concentration}{art:essential-labor:s-essential-goods-concentration}
The prohibition on dominant market positions under \ref{art:worker-coop-sector:s-cooperative-markets} extends explicitly to any cooperative or federation of cooperatives that withholds, stockpiles, or restricts the supply of a staple food item under \ref{art:essential-labor:s-food-stability-implementation} for the purpose of affecting its price or availability. The \gls{gls:financial-regulatory-body} monitors staple food supply chains for this purpose on an ongoing basis, in addition to responding to complaints. This Section governs enterprise-level conduct only; nothing in this Article establishes rationing or purchase limits on individual residents, whose entitlement under \ref{art:essential-labor:s-food-stability-implementation} remains unconditional regardless of market conditions.

\Article{Regulatory Commons Bodies}{art:regulatory-bodies}

\ArtSection{Definition and Purpose}{art:regulatory-bodies:s-definition-purpose}
Regulatory Commons Bodies are a distinct institutional category exercising independent regulatory authority within legislatively defined mandates. They combine democratic legitimacy, technical competence, and affected community representation to set and enforce standards that protect Commons assets, public health, worker safety, and ecological limits.

\ArtSection{Tripartite Composition}{art:regulatory-bodies:s-tripartite-composition}
Each Body has fifteen members, composed of three equal streams of five. The Democratic Stream consists of persons selected by sortition from the general civic membership with no financial interest in the regulated sector. The Expert Stream consists of persons selected by sortition from a domain-specific qualification pool assessed by the \gls{gls:qualification-commons} under Article \ref{art:sortition-design}, \ref{art:sortition-design:s-qualification-commons}. The Affected Community Stream consists of persons selected by sortition from communities most directly affected by the regulatory domain. The \gls{gls:sortition-legislature} may set a larger total for a specific Body by supermajority vote, provided the three streams remain equal in size.

\ArtSection{Sector Exclusion}{art:regulatory-bodies:s-sector-exclusion}
Persons who hold an elevated or non-standard stake in a specific regulatory outcome -- through executive or managerial roles, federation leadership positions, a currently pending or recently adjudicated matter before that Body involving their own cooperative, or income from outside consulting, advisory, or contracted work performed for entities in the regulated sector -- are excluded from all three streams. Ordinary worker-ownership under Article \ref{art:worker-coop-sector}, held on equal terms with other members of a cooperative in the regulated sector, does not by itself disqualify.

The time period defining a "recently adjudicated" matter under this Section, and the threshold defining disqualifying outside consulting, advisory, or contracted income, are set by the \gls{gls:sortition-legislature}, reviewed annually, independent of the five-year membership term structure of \ref{art:regulatory-bodies:s-terms-transparency}, and published through the \gls{gls:transparency-engine}.

\ArtSection{Mandate and Independence}{art:regulatory-bodies:s-mandate-independence}
The \gls{gls:sortition-legislature} defines each Body's mandate. Within that mandate, the Body acts independently. Regulated entities may appeal regulatory decisions only to the \gls{gls:constitutional-court} on the specific constitutional question of whether the Body exceeded its mandate -- not on policy grounds.

\ArtSection{Terms and Transparency}{art:regulatory-bodies:s-terms-transparency}
Members serve five-year staggered non-renewable terms. All deliberations, evidence considered, dissenting opinions, enforcement actions, and communications with regulated entities are logged in real time by the \gls{gls:transparency-engine}.

\ArtSection{Major Regulatory Domains}{art:regulatory-bodies:s-major-regulatory-domains}
The following Regulatory Commons Bodies are constitutionally mandated:
\begin{itemize}
    \item The \textbf{\gls{gls:environmental-regulatory-body}}, operating under Article \ref{art:ecological-limits}, sets and enforces emissions standards, water quality standards, land use requirements, and extraction limits. It monitors and audits enterprises regularly to ensure all standards are being met, and submits investigation findings to the \gls{gls:security-enforcement-review-panel} under Article \ref{art:internal-security}, \ref{art:internal-security:s-security-enforcement-review} for authorization of specific enforcement actions.

    \item The \textbf{\gls{gls:public-health-safety-regulatory-body}} oversees pharmaceutical safety, medical device safety, food safety, and public health standards.

    \item The \textbf{\gls{gls:worker-safety-regulatory-body}} sets minimum workplace safety floors below which cooperative democratic decisions may not fall.

    \item The \textbf{\gls{gls:financial-regulatory-body}} enforces the prohibitions enumerated in Article \ref{art:finance-money}, monitors financial flows within the cooperative and Commons sectors for violations, and classifies new financial instruments as permitted or provisionally prohibited. It also monitors cooperative market structures for the emergence of dominant positions that enable power over other cooperatives or over consumers' access to essential goods, as prohibited by Article \ref{art:worker-coop-sector}, \ref{art:worker-coop-sector:s-cooperative-markets}. It may order structural remedies including cooperative disaggregation, mandatory supply access, and federation restructuring. Investigation findings requiring enforcement action are submitted to the \gls{gls:security-enforcement-review-panel} under Article \ref{art:internal-security}, \ref{art:internal-security:s-security-enforcement-review}.

    \item The \textbf{\gls{gls:national-cooperative-standards-body}} sets and enforces the minimum size threshold at which an enterprise must organize as a worker cooperative under Article \ref{art:worker-coop-sector}, \ref{art:worker-coop-sector:s-cooperative-mandate}, including determinations of commonly-controlled or contractually-bound enterprise clusters for threshold purposes under that Section. It coordinates with the \gls{gls:financial-regulatory-body} on cases where enterprise-structuring findings are also relevant to an assessment of dominant market position under Article \ref{art:worker-coop-sector}, \ref{art:worker-coop-sector:s-cooperative-markets}, but its own mandate is limited to pre-formation classification and does not extend to post-formation market conduct.
\end{itemize}

\ArtSection{Financial Regulation}{art:regulatory-bodies:s-financial-regulation}
Provisional prohibitions referenced in \ref{art:regulatory-bodies:s-major-regulatory-domains} are referred to the \gls{gls:sortition-legislature} within thirty days of classification. If the Legislature does not act within ninety days of referral, the provisional prohibition becomes permanent. The Legislature may lift any prohibition by simple majority vote. The \gls{gls:financial-regulatory-body} may not itself create new categories of permitted financial activity -- that authority rests with the Legislature.

Additional Regulatory Bodies may be created by the \gls{gls:sortition-legislature} by supermajority vote. Constitutionally mandated Bodies may not be dissolved without a constitutional amendment under Article \ref{art:ordinary-amend}.